% ==== Cuts from pass 1: Appendix Data and code availability ====
% ---- from: Data and code availability (notes block) ----
% ---- NOTES: Data and code availability ----
% Decision for you: the author contributions now appear twice, in the Statement of Originality and in the last paragraph below. Keep both and they must stay word for word the same; keep one and point to it from the other. The two differ in one sentence now: "All authors read and approved the final manuscript" is in your paper and below, but not in the Statement of Originality.
% Decision for you: whether the contributions paragraph below also says A.P. wrote the QM9 pipeline. Your paper names only NoiseInject, so the sentence below now matches the paper. Every commit in the qsar_qm_models repository is yours, so adding it would be accurate.
% Decision for you: the licence for KIRBy. Your earlier plan was MIT or Apache-2.0. The repository has no licence file and is private, so the KIRBy block has placeholders, and the last sentence of the second data paragraph holds only once it is public.
% Decision for you: whether to give each public dataset's licence, as a column in the dataset table of Appendix~\ref{app:datasets}. Now the text gives licences only for OpenADMET, ChEMBL and the Biogen set; QM9 has none stated. The reaction-yield, antibody, nucleotide and gene-expression licences depend on which copy was downloaded and were not checked.
% Decision for you: the Python version for KIRBy. The package metadata declares 3.8 or later. The repository's setup guide gives 3.9 as the minimum and recommends 3.10 to 3.11. The university cluster has 3.9. The software block below says 3.8 or later, as the package metadata declares, until you pick one.
% Decision for you: the licence of the QM9 pipeline release. The qsar_qm_models repository holds a GNU GPL v3 licence file, so the block below says GPL v3. Your earlier plan was MIT or Apache-2.0. Keep GPL v3, or change it; the licence line and the "restrictions to use by non-academics" line both follow from your choice.
% Decision for you: your NoiseInject paragraph (restored as you wrote it) says "six regression and six classification noise injection strategies". Chapter~\ref{chap:robustness} and Section~\ref{sec:noise_injection} now define seven noise conditions. Keep the package's own count, or change it to match the chapters.
% ---- END NOTES ----
% ---- from: Data and code availability (TODO(you) lines above the software blocks and the Lhasa paragraph; now in DECISIONS.md) ----
% TODO(you): the DOI archives v0.3.0 (3 June 2026), not 1.0.0; replace it once a 1.0.0 Zenodo release exists.
% TODO(you): no release or Zenodo DOI yet; project name is the repository name until a release name is chosen. Confirm the operating systems it has been built on; Linux and macOS are assumed, not checked.
% TODO(you): make the repository public, confirm the URL, and add the Zenodo DOI and licence before deposit.
% TODO(you): confirm with Lhasa that the curated hERG sets cannot be redistributed, the wording they want, and whether FLuID can be listed as available on request.
